\documentclass[11pt,a4paper]{article}

\usepackage[utf8]{inputenc}
\usepackage[T1]{fontenc}
\usepackage{lmodern}
\usepackage{amsmath, amssymb, amsfonts}
\usepackage{bm}
\usepackage{geometry}
\usepackage{graphicx}
\usepackage{hyperref}
\usepackage{authblk}
\usepackage{cite}
\usepackage{physics}
\usepackage{mathtools}

\geometry{margin=2.4cm}

\title{\textbf{RDCC and Euclid: A Survey-Ready Framework for Testing Relaxation-Driven Cosmology}}
\author{Michael Lehmann \\ \texttt{mi.lehmann@gmx.de}}
\date{April 2026}

\begin{document}
\maketitle

\begin{abstract}
The Relaxation-Driven Cyclic Cosmology (RDCC) introduces a scale-dependent 
relaxation parameter $\alpha$ that modifies the gravitational potential and 
induces distinct geometric and topological signatures in weak-lensing 
observables. This paper provides a Euclid-focused overview of RDCC, summarizing 
its theoretical foundations, its observational signatures, and its compatibility 
with Euclid's survey design, data products, and analysis pipelines. We present 
Euclid-level forecasts, identify optimal RDCC probes, and outline a practical 
integration strategy for testing RDCC within the Euclid weak-lensing and 
large-scale structure programs. This document is intended as a concise, 
survey-oriented reference for Euclid collaborators and theorists.
\end{abstract}
\section{Introduction}

The Euclid mission provides an unprecedented opportunity to test extensions of 
$\Lambda$CDM through high-precision weak-lensing and galaxy-clustering 
measurements. The Relaxation-Driven Cyclic Cosmology (RDCC) is a theoretically 
motivated alternative that modifies the gravitational potential via a 
scale-dependent relaxation parameter $\alpha$. This modification induces 
geometric, morphological, and topological signatures in the convergence field 
that are directly accessible to Euclid's weak-lensing pipeline.

This paper summarizes the RDCC framework from a Euclid perspective. We focus on:
(i) the theoretical structure of RDCC, (ii) its observational signatures in 
Euclid data products, (iii) Euclid-level forecasts for $\alpha$, and 
(iv) practical integration pathways for testing RDCC within Euclid's analysis 
infrastructure.

\section*{Nomenclature}

\begin{tabular}{ll}
$\alpha$ & RDCC relaxation parameter; controls scale-dependent modifications of the potential \\[4pt]

$\Phi$ & Gravitational potential \\[4pt]

$\Phi_{\mathrm{RDCC}}$ & Modified potential in the RDCC model \\[4pt]

$\Phi_{\Lambda\mathrm{CDM}}$ & Standard potential in the $\Lambda$CDM model \\[4pt]

$\mathcal{R}(k,z;\alpha)$ & RDCC relaxation kernel; scale- and redshift-dependent modification \\[4pt]

$k$ & Wavenumber in Fourier space \\[4pt]

$z$ & Redshift \\[4pt]

$\kappa$ & Weak-lensing convergence field \\[4pt]

$\gamma$ & Shear field \\[4pt]

$P_{\kappa}(l)$ & Convergence power spectrum as a function of multipole $l$ \\[4pt]

$\mathrm{MF}_{i}$ & Minkowski functionals (area, boundary length, Euler characteristic) \\[4pt]

$\mathcal{D}$ & Persistence diagram in persistent homology \\[4pt]

$W_{2}$ & Quadratic optimal-transport distance between persistence diagrams \\[4pt]

$\sigma(\alpha)$ & Forecasted uncertainty on the RDCC parameter $\alpha$ \\[4pt]

$\Delta \log Z$ & Evidence contrast between RDCC and $\Lambda$CDM \\[4pt]

$v_{k}$ & Fisher channel (eigenvector of the information-tensor decomposition) \\[4pt]

$\lambda_{k}$ & Eigenvalue of the $k$-th Fisher channel (information contribution) \\[4pt]
\end{tabular}


\section{RDCC: Theoretical Overview}

RDCC modifies the gravitational potential through a relaxation kernel


\[
    \Phi_{\mathrm{RDCC}}(k,z)
    =
    \Phi_{\Lambda\mathrm{CDM}}(k,z)
    \cdot
    \mathcal{R}(k,z;\alpha),
\]


where $\mathcal{R}$ introduces scale-dependent deviations that become relevant 
at $k \gtrsim 0.1\,h\,\mathrm{Mpc}^{-1}$. These deviations propagate into the 
lensing potential, affecting:

\begin{itemize}
    \item the convergence field $\kappa$,
    \item shear correlations,
    \item peak and void statistics,
    \item Minkowski functionals,
    \item persistent homology and transport-based topology.
\end{itemize}

RDCC is fully compatible with Euclid's tomographic binning and angular 
resolution, and its signatures are strongest at intermediate multipoles 
accessible to Euclid's VIS instrument.
\section{RDCC Signatures in Euclid Observables}

RDCC induces characteristic distortions in Euclid weak-lensing data:

\subsection*{4.1 Power spectra}
Moderate-scale suppression or enhancement depending on $\alpha$; partially 
degenerate with $\Lambda$CDM parameters.

\subsection*{4.2 Peak and void statistics}
RDCC shifts peak heights and void depths, providing strong sensitivity to 
relaxation-driven potential changes.

\subsection*{4.3 Minkowski functionals}
Morphological distortions in area, boundary length, and Euler characteristic.

\subsection*{4.4 Persistent homology}
RDCC modifies the birth--death structure of topological features, producing 
robust, noise-resistant signatures.

\subsection*{4.5 Transport-based topology}
Optimal-transport distances between persistence diagrams provide a 
high-sensitivity probe of RDCC-induced non-Gaussian structure.
\section{Euclid-Level Forecasts}

Using Euclid-like noise, resolution, and tomographic binning, RDCC forecasts 
yield:


\[
    \sigma(\alpha) \sim \mathcal{O}(10^{-2}),
\]


with strongest constraints from:

\begin{itemize}
    \item transport-based topological statistics,
    \item peak and void counts,
    \item neural-compressed summaries,
    \item combined multi-probe likelihoods.
\end{itemize}

Fisher-channel analysis shows that $\sim 70\%$ of RDCC information resides in 
the leading channel, enabling efficient compression and robust inference.
\section{Degeneracy Structure and Euclid Synergy}

RDCC variations can partially mimic $\Lambda$CDM parameter shifts. Euclid's 
multi-probe design helps break these degeneracies:

\begin{itemize}
    \item tomographic binning separates redshift-dependent relaxation effects,
    \item cross-correlation with galaxy clustering reduces null directions,
    \item topological probes provide strong orthogonality to $\Lambda$CDM-aligned manifolds.
\end{itemize}
\section{Integration into Euclid Pipelines}

RDCC can be tested within Euclid's existing analysis framework:

\begin{itemize}
    \item \textbf{Forward modeling:} RDCC can be incorporated into lensing 
    simulations via modified potential kernels.

    \item \textbf{Summary statistics:} RDCC signatures are compatible with 
    Euclid's power-spectrum pipeline and can be extended to include peaks, 
    Minkowski functionals, and topology.

    \item \textbf{Likelihood integration:} RDCC likelihood components can be 
    added modularly to Euclid's cosmological parameter estimation framework.

    \item \textbf{Validation:} The RDCC Benchmark Suite provides 
    simulation-based calibration and cross-probe consistency tests.
\end{itemize}
\section{Conclusion}

RDCC provides a theoretically motivated and observationally testable extension 
of $\Lambda$CDM. Its geometric and topological signatures are well matched to 
Euclid's weak-lensing capabilities, and Euclid-level forecasts demonstrate that 
RDCC can be constrained with high precision. This paper outlines a practical 
pathway for integrating RDCC into Euclid's analysis pipelines and provides a 
survey-ready framework for testing relaxation-driven cosmology.
\end{document}